\chapter{INTRODUCTION}
\section{Background}
The advent of Large Language Models (LLMs) has fundamentally transformed the software development process. These models have become an integral part of modern workflows, assisting developers in tasks such as code generation, debugging, and documentation. Despite their success, studies such as \cite{Kaplan2020Scaling} show that LLMs rely on massive datasets and substantial computational resources to achieve improved performance. This dependence on high-performance Graphics Processing Units (GPUs) for training and inference makes them difficult to deploy and run locally in resource-constrained environments.
\section{Motivation}
Currently, state-of-the-art LLMs are primarily accessed through APIs for code generation. While smaller open-source models like Llama-2 7B \cite{Touvron2023Llama2} are available to run on single consumer GPU, their performance often falls short due to limited parameter scale. In our initial research, we found that generating Python code locally using these smaller base models resulted in poor code quality and high error rates. To bridge this gap, we explored fine-tuning, a process that specializes general models for specific tasks \cite{Howard2018ULMFiT}. However, full fine-tuning is computationally prohibitive for students under resource constraints. This can be addressed using Parameter-Efficient Fine-Tuning (PEFT) through Quantized Low-Rank Adaptation (QLoRA) \cite{Dettmers2023}, which quantizes the model to 4-bit precision and trains a minimal set of adapter weights \cite{Hu2021}, achieving performance comparable to full fine-tuning on limited hardware. Consequently, our goal is to fine-tune Llama-2 7B using QLoRA on a high-quality Python dataset, which we will prepare (manually and with the help of LLMs).

\section{Problem Definition}
The standard Llama-2 7B model suffers from a significant quality gap, frequently producing syntactically incorrect or logically flawed Python code compared to massive proprietary systems \cite{Diehl2024Llama2Benchmark}. Furthermore, resource constraints make full fine-tuning inaccessible, as it exceeds the VRAM capacity of consumer-grade GPUs like the NVIDIA T4 or RTX 3060. It also faces persistent logic issues, often guessing code structures and hallucinating during complex programming tasks rather than following a planned approach. 

\section{Objectives}
The primary objective of this project is:
\begin{itemize}
    \item To use QLoRA for compression and fine tuning of Llama-2 7B model so that it can run on standard consumer hardware for python code generation.
\end{itemize}

\section{Scope and Applications}
This project focuses on specializing the Llama-2 7B model for Python programming, specifically improving its ability to handle complex logic and various libraries. By using 4-bit QLoRA, the training is kept within a strict 14GB VRAM limit, ensuring the system remains compatible with consumer-level hardware and Google Colab. The model is trained using dataset either gathered manually or sourced from datasets like Alpaca-18k, to teach it problem-solving. Success will be measured by comparing the original and fine-tuned models against a series of Python coding tasks.

\section*{Potential Applications}
This project offers various practical benefits:
\begin{itemize}
    \item \textbf{IDE Integration:} Enables local integration within development environments to provide real-time, offline coding assistance.
    
    \item \textbf{Data Privacy:} Ensures that sensitive code and data are processed entirely on-site, preserving user privacy.
    
    \item \textbf{AI-Powered Debugging Assistant:} Assists in identifying errors and suggesting fixes and improvements in Python code.
    
    \item \textbf{Code Refactoring and Optimization:} Enhances existing code by removing redundancy and suggesting more efficient implementations.
    
    \item \textbf{Code Generation for Competitive Programming:} Assists students in solving algorithmic problems from platforms such as LeetCode and Codeforces.

    \item \textbf{Educational Assistant:} Supports students in learning programming concepts by generating example solutions and explanations.
\end{itemize}